\documentclass[12pt]{article}
\usepackage{amsmath, amssymb, amsfonts}
\usepackage{graphicx}
\usepackage{cite}
\usepackage{geometry}
\usepackage{hyperref}
\geometry{margin=1in}

\title{\textbf{Emergent Geometry from Entanglement:\\
A Numerical Demonstration using Heisenberg Spin Chains}}

\author{Paul Jarvis\\
Independent Researcher, UK\\
\texttt{mrpaulwjarvis@gmail.com}\\
ORCID: 0009-0009-8933-857X}

\date{\today}

\begin{document}

\maketitle

\begin{abstract}
We provide a model-independent numerical demonstration of the emergence of smooth spatial geometry directly from the entanglement structure of a pure quantum state. Utilizing exact diagonalization of an isotropic antiferromagnetic Heisenberg spin chain, we compute the ground-state mutual information matrix without presupposing a background manifold or metric framework. By normalizing these raw quantum correlations, we construct a discrete adjacency matrix $A_{ij}$ and its corresponding graph Laplacian $L$. Spectral analysis reveals that a continuous one-dimensional Laplace--Beltrami operator emerges naturally in the continuum limit, characterized by a constant zero-mode eigenvector and a quadratic scaling of the lowest non-zero eigenvalues $\lambda_k \propto k^2$. Furthermore, we show that the power-law decay of mutual information in this gapless critical state organically enforces Neumann-type boundary conditions without fine-tuning, establishing substrate criticality as a necessary condition for smooth emergent spatial structures.
\end{abstract}

\section{Introduction}

A foundational paradigm in modern high-energy physics is that spacetime geometry is not an absolute background container, but an emergent macro-property woven from the entanglement structure of quantum states. While this concept has been explored extensively through holographic dualities and quantum error-correcting codes, explicit numerical verifications on simple many-body systems remain sparse.

This work provides an operational, data-driven protocol to extract smooth spatial differential operators directly from a pure quantum ground state. Rather than treating geometry as an input, we treat the ground state of a condensed matter system---the isotropic antiferromagnetic Heisenberg spin chain---as a raw relational data set. By computing the mutual information matrix, we build a discrete relational graph and analyze its spectral properties.

Our results demonstrate that a true spatial continuum limit emerges organically, indicating that the differential structure of space can arise as a mathematical consequence of structured quantum correlations.

\section{Methodology and Operator Construction}

We consider a one-dimensional, open isotropic antiferromagnetic Heisenberg spin chain governed by the Hamiltonian
\begin{equation}
H = J \sum_{i=1}^{N-1} \mathbf{S}_i \cdot \mathbf{S}_{i+1},
\end{equation}
where $J > 0$ sets the antiferromagnetic coupling and $\mathbf{S}_i = (\sigma_x, \sigma_y, \sigma_z)/2$ are spin-$1/2$ operators. The exact ground state $|\psi_0\rangle$ is obtained via exact diagonalization for system sizes $N = 6$ to $N = 12$.

For each pair of sites $(i,j)$, we compute the reduced density matrix $\rho_{ij}$ and the von Neumann mutual information
\begin{equation}
I(i,j) = S(\rho_i) + S(\rho_j) - S(\rho_{ij}),
\end{equation}
where $S(\rho) = -\mathrm{Tr}(\rho \ln \rho)$.

To construct a scale-independent relational graph, we define the normalized adjacency matrix
\begin{equation}
A_{ij} = \frac{I(i,j)}{\max_{k,\ell} I(k,\ell)}, \qquad A_{ii} = 0.
\end{equation}

The corresponding graph Laplacian is
\begin{equation}
L = D - A,
\end{equation}
where $D$ is the degree matrix with $D_{ii} = \sum_j A_{ij}$.

\section{Spectral Results and Continuum Verification}

We perform an exact eigenvalue decomposition of $L$ for all system sizes considered. The results exhibit three definitive signatures of emergent one-dimensional geometry.

\subsection{Constant Zero-Mode}

A Laplace--Beltrami operator on a one-dimensional manifold with Neumann boundaries must possess a zero eigenvalue with a constant eigenvector. Numerically, we obtain


\[
\lambda_0 < 10^{-15}, \qquad v_0 = (1,1,\ldots,1)^T,
\]


confirming that the information-derived operator preserves the constant mode.

\subsection{Emergent Neumann Boundary Conditions}

Higher eigenmodes exhibit smooth, cosine-like profiles. Because the mutual information decays smoothly toward the boundaries of the open chain, the operator naturally enforces zero spatial derivative at the endpoints. Neumann-type boundary conditions therefore arise organically from the correlation structure.

\begin{figure}[h]
\centering
\includegraphics[width=0.75\textwidth]{eigenmodes.png}
\caption{Representative eigenmodes of the graph Laplacian for $N=12$, showing excellent agreement with cosine functions satisfying Neumann boundary conditions.}
\label{fig:eigenmodes}
\end{figure}

\subsection{Quadratic Dispersion Scaling}

For a smooth one-dimensional manifold, the low-lying eigenvalues must scale as $\lambda_k \propto k^2$. To isolate the continuum-relevant region, we restrict attention to the first four non-zero eigenvalues. A quadratic fit yields $R^2 > 0.998$ for all system sizes.

\begin{figure}[h]
\centering
\includegraphics[width=0.75\textwidth]{spectrum.png}
\caption{Low-lying Laplacian eigenvalues for $N=12$ with quadratic fit. Only the first four non-zero modes are shown, corresponding to the continuum regime.}
\label{fig:spectrum}
\end{figure}

Finite-size scaling of the lowest non-zero eigenvalue further confirms the continuum limit:


\[
\lambda_1(N) \propto \frac{1}{N^2}.
\]



\begin{figure}[h]
\centering
\includegraphics[width=0.75\textwidth]{scaling_lambda1.png}
\caption{Finite-size scaling of the lowest non-zero eigenvalue $\lambda_1$ across $N = 6, 8, 10, 12$. The linear dependence on $1/N^2$ confirms the expected continuum scaling.}
\label{fig:scaling}
\end{figure}

\section{The Role of Substrate Criticality}

The isotropic Heisenberg chain is a gapless system described by a conformal field theory, leading to power-law decay of mutual information. If the system is perturbed into a gapped phase, correlations decay exponentially and the smooth sinusoidal eigenmodes collapse into localized clusters. Numerical tests confirm that only critical, power-law-correlated substrates support smooth emergent geometry. Substrate criticality is therefore a necessary condition for the emergence of continuous spatial structure.

\begin{figure}[h]
\centering
\includegraphics[width=0.75\textwidth]{mi_heatmap.png}
\caption{Mutual-information matrix for $N=12$, showing smooth power-law decay and boundary-induced symmetry consistent with Neumann conditions.}
\label{fig:mi}
\end{figure}

\section{Conclusion}

We have presented a numerical demonstration that continuous spatial manifolds and their associated differential operators can emerge directly from the entanglement structure of a pure quantum state. By analyzing the spectral properties of a graph Laplacian constructed from normalized mutual information, we verified the appearance of a constant zero mode, Neumann boundary conditions, and quadratic dispersion scaling. These results provide a concrete, data-driven validation tool for information-theoretic approaches to emergent geometry.

\section*{Data Availability}

All numerical data and code used in this study are available from the author upon request.

\bibliographystyle{unsrt}
\begin{thebibliography}{9}

\bibitem{raamsdonk}
M.~Van Raamsdonk,
\textit{Building Spacetime from Entanglement},
Gen. Rel. Grav. \textbf{42}, 2323--2329 (2010).

\bibitem{maldacena}
J.~Maldacena and L.~Susskind,
\textit{Cool Horizons for Entangled Black Holes},
Fortsch. Phys. \textbf{61}, 781--811 (2013).

\bibitem{bethe}
H.~Bethe,
\textit{On the Theory of Metals. I. Eigenvalues and Eigenfunctions of the Linear Atomic Chain},
Z. Phys. \textbf{71}, 205--226 (1931).

\bibitem{amico}
L.~Amico, R.~Fazio, A.~Osterloh, and V.~Vedral,
\textit{Entanglement in Many-Body Systems},
Rev. Mod. Phys. \textbf{80}, 517--576 (2008).

\end{thebibliography}

\end{document}
